%% Generated by Sphinx.
\def\sphinxdocclass{report}
\documentclass[a4paper,11pt,oneside,english,openany]{sphinxmanual}
\ifdefined\pdfpxdimen
   \let\sphinxpxdimen\pdfpxdimen\else\newdimen\sphinxpxdimen
\fi \sphinxpxdimen=.75bp\relax
\ifdefined\pdfimageresolution
    \pdfimageresolution= \numexpr \dimexpr1in\relax/\sphinxpxdimen\relax
\fi
%% let collapsible pdf bookmarks panel have high depth per default
\PassOptionsToPackage{bookmarksdepth=5}{hyperref}

\PassOptionsToPackage{booktabs}{sphinx}
\PassOptionsToPackage{colorrows}{sphinx}

\PassOptionsToPackage{warn}{textcomp}
\usepackage[utf8]{inputenc}
\ifdefined\DeclareUnicodeCharacter
% support both utf8 and utf8x syntaxes
  \ifdefined\DeclareUnicodeCharacterAsOptional
    \def\sphinxDUC#1{\DeclareUnicodeCharacter{"#1}}
  \else
    \let\sphinxDUC\DeclareUnicodeCharacter
  \fi
  \sphinxDUC{00A0}{\nobreakspace}
  \sphinxDUC{2500}{\sphinxunichar{2500}}
  \sphinxDUC{2502}{\sphinxunichar{2502}}
  \sphinxDUC{2514}{\sphinxunichar{2514}}
  \sphinxDUC{251C}{\sphinxunichar{251C}}
  \sphinxDUC{2572}{\textbackslash}
\fi
\usepackage{cmap}
\usepackage[T1]{fontenc}
\usepackage{amsmath,amssymb,amstext}
\usepackage{babel}





\usepackage[Bjarne]{fncychap}
\usepackage{sphinx}

\fvset{fontsize=auto}
\usepackage{geometry}


% Include hyperref last.
\usepackage{hyperref}
% Fix anchor placement for figures with captions.
\usepackage{hypcap}% it must be loaded after hyperref.
% Set up styles of URL: it should be placed after hyperref.
\urlstyle{same}

\addto\captionsenglish{\renewcommand{\contentsname}{Contents:}}

\usepackage{sphinxmessages}
\setcounter{tocdepth}{1}


		\renewcommand{\familydefault}{\rmdefault}
        \usepackage{setspace} % Required for more spacing options
		\parindent 0px
		\setstretch{1.1}
        \usepackage[T1]{fontenc}
        \usepackage{charter}
        \usepackage{amsmath}
        \usepackage{amsfonts}
        \usepackage{amssymb}
        \usepackage{fancyhdr}
        \usepackage{arydshln}
        \usepackage{xcolor}
        \usepackage{makeidx}  % To create an index
        \usepackage{tocloft}  % To customize Table of Contents

        % Fancy header and footer
        \pagestyle{fancy}
        \fancyhead[L]{\leftmark}
        \fancyhead[R]{\thepage}
        \fancyfoot[C]{}

        % Table of Contents: Formatting (removing subsections)
        \setcounter{tocdepth}{1}  % Limit ToC to sections only (remove subsections)
        \renewcommand{\cftsecfont}{\normalfont\bfseries}  % Bold section titles
        \renewcommand{\cftsecindent}{0em}                 % No indentation for sections
        \setlength{\cftbeforesecskip}{0.5ex}              % Space before section titles

        % Index Configuration (simplified to sections only)
        \makeindex  % Generate index

        % Page layout adjustments
		\usepackage{geometry}
        \geometry{top=1.5cm, bottom=1.5cm, left=1.5cm, right=1.5cm}
        \usepackage{hyperref}
        \hypersetup{colorlinks=true,linkcolor=violet,filecolor=magenta,urlcolor=blue}

        % Better verbatim environment for code blocks
        %\usepackage{fvextra}
        \usepackage{tcolorbox}
        
        \newtcbox{\tbox}{colframe=black,colback=gray!5,boxrule=0.7pt,arc=4pt,boxsep=5pt}
    

\title{TGE}
\date{May 20, 2026}
\release{1.0.0}
\author{AE,SB,BB}
\newcommand{\sphinxlogo}{\vbox{}}
\renewcommand{\releasename}{Release}
\makeindex
\begin{document}

\ifdefined\shorthandoff
  \ifnum\catcode`\=\string=\active\shorthandoff{=}\fi
  \ifnum\catcode`\"=\active\shorthandoff{"}\fi
\fi

\pagestyle{empty}
\sphinxmaketitle
\pagestyle{plain}
\sphinxtableofcontents
\pagestyle{normal}
\phantomsection\label{\detokenize{index::doc}}
\sphinxstepscope



\sphinxstepscope


\chapter{grid module}
\label{\detokenize{grid:module-grid}}\label{\detokenize{grid:grid-module}}\label{\detokenize{grid::doc}}\index{module@\spxentry{module}!grid@\spxentry{grid}}\index{grid@\spxentry{grid}!module@\spxentry{module}}

\section{Tapered Gridded Estimator (TGE)}
\label{\detokenize{grid:tapered-gridded-estimator-tge}}
\sphinxAtStartPar
For a given \sphinxstylestrong{input uv\sphinxhyphen{}fits file}, we taper the sky response of the telescope by convolving visibilities with a function \(\tilde{w}(\vec{\textbf{U}})=\pi \theta_w^2 \textbf{exp}[-\pi^2 \theta_w^2 \textbf{U}^2]\) (which is the Fourier Transform of a window function \(\mathcal{W}(\vec{\boldsymbol{\theta}}) = \textbf{exp}[-\theta^2/\theta_w^2]\) which falls off to a value close to zero well before the first null of the telescope’s primary beam pattern ) and gridded on a square grid in the \(uv\)\sphinxhyphen{}plane using,
\begin{equation*}
\begin{split}\begin{align}
\boxed{\mathcal{V}_{cg}^{P}(\vec{\textbf{U}}_g,\nu_a)=\sum_{i}\  \tilde{w}(\vec{\textbf{U}}_g - \vec{\textbf{U}}_i)\ \mathcal{V}^{P}(\vec{\textbf{U}}_i,\nu_a)\ F_i^{P}(\nu_a)}
\end{align}\end{split}
\end{equation*}
\sphinxAtStartPar
The sum is over all the baselines. Here \(\nu_a\) refers to the different frequency channels, \(\vec{\textbf{U}}_i\) is \(i^{\rm th}\) baseline. \(\mathcal{V}_{cg}^{P}\) refers to the convolved visibility at grid point \(g\) and by the superscript \(P\) we denote the polarization. \(\vec{\textbf{U}}_g\) is the baseline corresponding to this grid point and \(F_i^{P}(\nu_a)\) incorporates the flagging information of the data corresponding to the frequency \(\nu_a\) and for the corresponding polarization \(P\). \(P\) can take be either in  \(\rm XX,YY,XY,YX\) basis or in the circular basis \(\rm RR,LL,RL,LR\), based on the telescope.

\sphinxAtStartPar
\(F_i^{P}(\nu_a)\) has a value ‘0’ if the data at a given baseline and frequency is flagged and ‘1’ otherwise. Note that the baselines \(\vec{\textbf{U}}_i\) here are defined at a fixed reference frequency \(\nu_c\) (\sphinxstylestrong{centered frequency of the fits file}), and in this analysis they are considered to be fixed as we vary the frequency channel \(\nu_a\) (although in principle they are not).

\sphinxAtStartPar
In the visibility data, we flag the data if the corresponding weight value is less than 0, otherwise we take it as it is.

\sphinxAtStartPar
We calculate convolved visibility \(\mathcal{V}_{cg}^{P}(\vec{\textbf{U}}_g,\nu_a)\) for a given no of frequency channels and for given polarizations of the data. A 4D array is returned, where \sphinxstylestrong{first one is the polarization axis}, next \sphinxstylestrong{two dimensions are values of convolved visibility at each grid points} (square grid) in \(uv\) plane and \sphinxstylestrong{last one is frequency axis}.

\clearpage

\sphinxAtStartPar
Here \(\theta_w = f\theta_0\), with \(\theta_0 = 0.6 \ \theta_{\rm FWHM}\). \(f\) is known as tapering parameter whose value can be suitably chosen. The effect of tapering is enhanced if the value of \(f\) is reduced. For instance \(f\) \textless{} 1 would suppress the sky response from the outer regions and the side\sphinxhyphen{}lobes of the PB pattern, whereas a large value \(f\) \textgreater{} 1 would imply very little tapering.
Preferably \(f\leq 1\) so that \(\mathcal{W}(\vec{\boldsymbol{\theta}})\) cuts off the sky response well before the first null of the primary beam.

\sphinxAtStartPar
The grid spacing in the \(uv\) plane is given by  \(\Delta{\text{U}}=\dfrac{\sqrt{ln2}}{2\pi\theta_{\rm eff}}\) with \(\theta_{\rm eff}=\dfrac{f\theta_0}{\sqrt{1+f^2}}\).

\sphinxAtStartPar
Meanwhile, those baselines are rejected for which \(|\vec{\textbf{U}}_i|<|\vec{\textbf{U}}|_{max}\). We effectively increase the size of visibility data twice ; original one and one for all the baselines with \(-\vec{\textbf{U}}_i\) as \(\mathcal{V}^{P}(-\vec{\textbf{U}}_i,\nu_a) = \left[\mathcal{V}^{P}(\vec{\textbf{U}}_i,\nu_a)\right]^{*}\).

\sphinxAtStartPar
For any baseline \(\vec{\textbf{U}}_i\), the contribution of visibilities to the grid points is taken for only those grid points for which are within(and on) a square grid of side length \(12\Delta{\text{U}}\) whose center is at grid point which is nearest to \(\vec{\textbf{U}}_i\) (The sides of this grid are along \(u\) and \(v\) axis). The weight function  \(\tilde{w}(\vec{\textbf{U}}_g - \vec{\textbf{U}}_i)\) falls considerably faster and we do not expect a significant contribution from the visibilities beyond this baseline separation.

\sphinxAtStartPar
\(\textcolor{blue}{\rm \textbf{Please note that in this analysis Baseline migration is not incorporated.}}\)

\clearpage


\section{Module Functions}
\label{\detokenize{grid:module-functions}}\index{mkbin() (in module grid)@\spxentry{mkbin()}\spxextra{in module grid}}

\begin{fulllineitems}
\phantomsection\label{\detokenize{grid:grid.mkbin}}
\pysigstartsignatures
\pysiglinewithargsret{\sphinxcode{\sphinxupquote{grid.}}\sphinxbfcode{\sphinxupquote{mkbin}}}{\sphinxparam{\DUrole{n}{GV}}\sphinxparamcomma \sphinxparam{\DUrole{n}{dU}}\sphinxparamcomma \sphinxparam{\DUrole{n}{Umax}}\sphinxparamcomma \sphinxparam{\DUrole{n}{FWHM}}\sphinxparamcomma \sphinxparam{\DUrole{n}{f}}\sphinxparamcomma \sphinxparam{\DUrole{n}{binUmax}}\sphinxparamcomma \sphinxparam{\DUrole{n}{binUmin}}\sphinxparamcomma \sphinxparam{\DUrole{n}{Nbin}}\sphinxparamcomma \sphinxparam{\DUrole{n}{Mg\_min}}\sphinxparamcomma \sphinxparam{\DUrole{n}{outf}}}{}
\pysigstopsignatures
\sphinxAtStartPar
Given a \sphinxstylestrong{UAPS convolved gridded visibility array} and suitable gridding parameters it identfies the grid points which are relevant and which grid points are in which bin.
Also the effective multipole value for a particular bin is given by :
\begin{equation*}
\begin{split}\ell_a = \frac{\sum_g\,w_g\,\ell_g}{\sum_g\,w_g}; \quad \text{with}\, w_g = M_g\end{split}
\end{equation*}
\sphinxAtStartPar
And \(\ell_g = 2\pi |\vec{\textbf{U}}|_{g}\)
\begin{quote}\begin{description}
\sphinxlineitem{Parameters}\begin{itemize}
\item {} 
\sphinxAtStartPar
\sphinxstyleliteralstrong{\sphinxupquote{GV}} (\sphinxstyleliteralemphasis{\sphinxupquote{np.ndarray}}) \textendash{} Convolved gridded visbility array for the \sphinxstylestrong{UAPS} visibilities. Different realizations have been put along frequency axis.

\item {} 
\sphinxAtStartPar
\sphinxstyleliteralstrong{\sphinxupquote{dU}} (\sphinxstyleliteralemphasis{\sphinxupquote{float}}) \textendash{} Grid spacing \(\Delta \text{U}\) in wavelength units.

\item {} 
\sphinxAtStartPar
\sphinxstyleliteralstrong{\sphinxupquote{Umax}} (\sphinxstyleliteralemphasis{\sphinxupquote{int}}\sphinxstyleliteralemphasis{\sphinxupquote{ or }}\sphinxstyleliteralemphasis{\sphinxupquote{float}}) \textendash{} The maximum value of baseline \(|\vec{\textbf{U}}|_{max}\). Baselines greater than this are rejected.

\item {} 
\sphinxAtStartPar
\sphinxstyleliteralstrong{\sphinxupquote{FWHM}} (\sphinxstyleliteralemphasis{\sphinxupquote{int}}\sphinxstyleliteralemphasis{\sphinxupquote{ or }}\sphinxstyleliteralemphasis{\sphinxupquote{float}}) \textendash{} The full width half maximum (FWHM) of telescope’s primary beam pattern.(in degrees)

\item {} 
\sphinxAtStartPar
\sphinxstyleliteralstrong{\sphinxupquote{f}} (\sphinxstyleliteralemphasis{\sphinxupquote{int}}\sphinxstyleliteralemphasis{\sphinxupquote{ or }}\sphinxstyleliteralemphasis{\sphinxupquote{float}}) \textendash{} Tapering parameter value.

\item {} 
\sphinxAtStartPar
\sphinxstyleliteralstrong{\sphinxupquote{binUmax}} (\sphinxstyleliteralemphasis{\sphinxupquote{int}}\sphinxstyleliteralemphasis{\sphinxupquote{ or }}\sphinxstyleliteralemphasis{\sphinxupquote{float}}) \textendash{} Maximum value of \(|\vec{\textbf{U}}|\) (in units of wavelength \(\lambda\)) to be used in the binning infomation. Grid points having baseline greater than this are rejected.

\item {} 
\sphinxAtStartPar
\sphinxstyleliteralstrong{\sphinxupquote{binUmin}} (\sphinxstyleliteralemphasis{\sphinxupquote{int}}\sphinxstyleliteralemphasis{\sphinxupquote{ or }}\sphinxstyleliteralemphasis{\sphinxupquote{float}}) \textendash{} Minimum value of \(|\vec{\textbf{U}}|\) (in units of wavelength \(\lambda\)) to be used in the binning infomation. Grid points having baseline lesser than this are rejected.

\item {} 
\sphinxAtStartPar
\sphinxstyleliteralstrong{\sphinxupquote{Nbin}} (\sphinxstyleliteralemphasis{\sphinxupquote{int}}) \textendash{} Number of annular bins the whole \(uv\) plane has been divided into.

\item {} 
\sphinxAtStartPar
\sphinxstyleliteralstrong{\sphinxupquote{Mg\_min}} (\sphinxstyleliteralemphasis{\sphinxupquote{int}}\sphinxstyleliteralemphasis{\sphinxupquote{ or }}\sphinxstyleliteralemphasis{\sphinxupquote{float}}) \textendash{} Minimum value of normalization factor \(M_g\) to be used. Grid points at which visibility self correlation is less than this are rejected.

\item {} 
\sphinxAtStartPar
\sphinxstyleliteralstrong{\sphinxupquote{outf}} (\sphinxstyleliteralemphasis{\sphinxupquote{string}}) \textendash{} The name of file in which the binning information is saved. This is saved in \sphinxstylestrong{.npz} format.

\end{itemize}

\end{description}\end{quote}
\subsubsection*{Notes}
\begin{itemize}
\item {} 
\sphinxAtStartPar
\(d\text{U}\) is the grid spacing in baseline units (\(\lambda\)).

\item {} 
\sphinxAtStartPar
\(N\!I\) is a 2D matrix, which contains the information which grid points are relevent. It contains value \(-2,-1,0,\cdots,N_{bin} - 1\). Grid points those satisfy \(NI \geq 0\) are relevant. The rest of them are rejected in our analysis.

\item {} 
\sphinxAtStartPar
\(ix\) is a 1D array, which contains the \(u\) index of grid points in terms of grid spacing \(d\text{U}\), i.e. \(ix_g = u_g/d\text{U}\). (only for relevant grid points).

\item {} 
\sphinxAtStartPar
\(iy\) is a 1D array, which contains the \(v\) index of grid points in terms of grid spacing \(d\text{U}\), i.e. \(iy_g = v_g/d\text{U}\). (only for relevant grid points).

\item {} 
\sphinxAtStartPar
\(ni\) is a 1D array, which contains information about which grid point falls into which annular bin. If the value is zero then it is residing in first bin, if 2 then in second bin and so on. Max value is \(N_{bin} - 1\). (only for relevant grid points).

\end{itemize}
\subsubsection*{Examples}

\begin{sphinxVerbatim}[commandchars=\\\{\}]
\PYG{g+gp}{\PYGZgt{}\PYGZgt{}\PYGZgt{} }\PYG{c+c1}{\PYGZsh{} BIN INFO}
\PYG{g+gp}{\PYGZgt{}\PYGZgt{}\PYGZgt{} }\PYG{k+kn}{import} \PYG{n+nn}{grid} \PYG{k}{as} \PYG{n+nn}{gd}
\PYG{g+go}{Imported Grid, Import Time : 19/05/26 | 07:48:35 PM}
\PYG{g+gp}{\PYGZgt{}\PYGZgt{}\PYGZgt{} }\PYG{c+c1}{\PYGZsh{} ===================================================}
\PYG{g+gp}{\PYGZgt{}\PYGZgt{}\PYGZgt{} }\PYG{c+c1}{\PYGZsh{} Set the gridding parameters in the TGE code}
\PYG{g+gp}{\PYGZgt{}\PYGZgt{}\PYGZgt{} }\PYG{n}{n1}      \PYG{o}{=} \PYG{l+m+mi}{0}     \PYG{c+c1}{\PYGZsh{} Starting channel number}
\PYG{g+gp}{\PYGZgt{}\PYGZgt{}\PYGZgt{} }\PYG{n}{n2}      \PYG{o}{=} \PYG{l+m+mi}{99}    \PYG{c+c1}{\PYGZsh{} End channel number}
\PYG{g+gp}{\PYGZgt{}\PYGZgt{}\PYGZgt{} }\PYG{n}{Flag}    \PYG{o}{=} \PYG{k+kc}{False} \PYG{c+c1}{\PYGZsh{} Apply actual flagging of data}
\PYG{g+gp}{\PYGZgt{}\PYGZgt{}\PYGZgt{} }\PYG{n}{Umax}    \PYG{o}{=} \PYG{l+m+mi}{250}   \PYG{c+c1}{\PYGZsh{} Baselines greater than this are rejected}
\PYG{g+gp}{\PYGZgt{}\PYGZgt{}\PYGZgt{} }\PYG{n}{FWHM}    \PYG{o}{=} \PYG{l+m+mi}{23}    \PYG{c+c1}{\PYGZsh{} Primary beam FWHM of the telescope in degrees}
\PYG{g+gp}{\PYGZgt{}\PYGZgt{}\PYGZgt{} }\PYG{n}{f}       \PYG{o}{=} \PYG{l+m+mf}{0.6}   \PYG{c+c1}{\PYGZsh{} Tapering parameter value}
\PYG{g+gp}{\PYGZgt{}\PYGZgt{}\PYGZgt{} }\PYG{n}{nstokes} \PYG{o}{=} \PYG{p}{[}\PYG{l+m+mi}{0}\PYG{p}{]}   \PYG{c+c1}{\PYGZsh{} 0 for XX and 1 for YY (polarizations to grid)}
\PYG{g+gp}{\PYGZgt{}\PYGZgt{}\PYGZgt{} }\PYG{c+c1}{\PYGZsh{} ===================================================}
\PYG{g+gp}{\PYGZgt{}\PYGZgt{}\PYGZgt{} }\PYG{n}{nrel}    \PYG{o}{=} \PYG{o}{\PYGZhy{}}\PYG{l+m+mi}{1}    \PYG{c+c1}{\PYGZsh{} grid along freq axis}
\PYG{g+gp}{\PYGZgt{}\PYGZgt{}\PYGZgt{} }\PYG{c+c1}{\PYGZsh{} For binning}
\PYG{g+gp}{\PYGZgt{}\PYGZgt{}\PYGZgt{} }\PYG{n}{Nbin}    \PYG{o}{=} \PYG{l+m+mi}{20}    \PYG{c+c1}{\PYGZsh{} No of annular bins the whole uv plane has been divided into}
\PYG{g+gp}{\PYGZgt{}\PYGZgt{}\PYGZgt{} }\PYG{n}{binUmin} \PYG{o}{=} \PYG{l+m+mf}{6.0}   \PYG{c+c1}{\PYGZsh{} min |U| used in getting binning information}
\PYG{g+gp}{\PYGZgt{}\PYGZgt{}\PYGZgt{} }\PYG{n}{binUmax} \PYG{o}{=} \PYG{l+m+mi}{220}   \PYG{c+c1}{\PYGZsh{} max |U| used in getting binning information}
\PYG{g+gp}{\PYGZgt{}\PYGZgt{}\PYGZgt{} }\PYG{n}{Mg\PYGZus{}min}  \PYG{o}{=} \PYG{l+m+mf}{0.01}  \PYG{c+c1}{\PYGZsh{} M\PYGZus{}g cutoff, grid points have less than this will be rejected}
\PYG{g+gp}{\PYGZgt{}\PYGZgt{}\PYGZgt{}}
\PYG{g+gp}{\PYGZgt{}\PYGZgt{}\PYGZgt{} }\PYG{n}{inpath}   \PYG{o}{=} \PYG{l+s+s1}{\PYGZsq{}}\PYG{l+s+s1}{/home/baijayanta/Complete\PYGZus{}Pipeline/uaps\PYGZus{}files/}\PYG{l+s+s1}{\PYGZsq{}}
\PYG{g+gp}{\PYGZgt{}\PYGZgt{}\PYGZgt{} }\PYG{n}{filename} \PYG{o}{=} \PYG{l+s+s1}{\PYGZsq{}}\PYG{l+s+s1}{1000007200\PYGZus{}out\PYGZus{}uaps\PYGZus{}15.fits}\PYG{l+s+s1}{\PYGZsq{}}
\PYG{g+gp}{\PYGZgt{}\PYGZgt{}\PYGZgt{} }\PYG{c+c1}{\PYGZsh{} name of input .fits file for which we will calculate the bin info}
\PYG{g+gp}{\PYGZgt{}\PYGZgt{}\PYGZgt{} }\PYG{n}{infile}   \PYG{o}{=} \PYG{n}{inpath} \PYG{o}{+} \PYG{n}{filename}
\PYG{g+gp}{\PYGZgt{}\PYGZgt{}\PYGZgt{} }\PYG{c+c1}{\PYGZsh{} grid it \PYGZsh{} ([Nstokes,Nu,Nv,Nc],(dU,Umax,FWHM,f))}
\PYG{g+gp}{\PYGZgt{}\PYGZgt{}\PYGZgt{} }\PYG{n}{GV}\PYG{p}{,} \PYG{n}{info} \PYG{o}{=} \PYG{n}{gd}\PYG{o}{.}\PYG{n}{grid}\PYG{p}{(}\PYG{n}{infile}\PYG{p}{,}\PYG{n}{n1}\PYG{p}{,}\PYG{n}{n2}\PYG{p}{,}\PYG{n}{nrel}\PYG{p}{,}\PYG{n}{Umax}\PYG{p}{,}\PYG{n}{FWHM}\PYG{p}{,}\PYG{n}{f}\PYG{p}{,}\PYG{n}{Flag}\PYG{p}{,} \PYG{n}{nstokes}\PYG{p}{)}
\PYG{g+go}{\PYGZlt{}\PYGZlt{}\PYGZlt{}\PYGZlt{}\PYGZlt{}\PYGZlt{}\PYGZlt{}\PYGZlt{}\PYGZlt{}\PYGZlt{}\PYGZlt{} TGE \PYGZgt{}\PYGZgt{}\PYGZgt{}\PYGZgt{}\PYGZgt{}\PYGZgt{}\PYGZgt{}\PYGZgt{}\PYGZgt{}\PYGZgt{}\PYGZgt{}\PYGZgt{}}
\PYG{g+go}{\PYGZlt{}\PYGZlt{} 19/05/26 | 07:48:35 PM \PYGZgt{}\PYGZgt{}}
\PYG{g+go}{Type  : Data}
\PYG{g+go}{\PYGZlt{}\PYGZlt{} 19/05/26 | 07:48:48 PM \PYGZgt{}\PYGZgt{}}
\PYG{g+go}{Elapsed   :   00:00:12.82}
\PYG{g+go}{\PYGZlt{}\PYGZlt{}\PYGZlt{}\PYGZlt{}\PYGZlt{}\PYGZlt{}\PYGZlt{}\PYGZlt{}\PYGZlt{}\PYGZlt{}\PYGZlt{} Done \PYGZgt{}\PYGZgt{}\PYGZgt{}\PYGZgt{}\PYGZgt{}\PYGZgt{}\PYGZgt{}\PYGZgt{}\PYGZgt{}\PYGZgt{}\PYGZgt{}}
\PYG{g+gp}{\PYGZgt{}\PYGZgt{}\PYGZgt{} }\PYG{c+c1}{\PYGZsh{} make bin info}
\PYG{g+gp}{\PYGZgt{}\PYGZgt{}\PYGZgt{} }\PYG{n}{gd}\PYG{o}{.}\PYG{n}{mkbin}\PYG{p}{(}\PYG{n}{GV}\PYG{p}{[}\PYG{l+m+mi}{0}\PYG{p}{]}\PYG{p}{,} \PYG{n}{info}\PYG{p}{[}\PYG{l+m+mi}{0}\PYG{p}{]}\PYG{p}{,} \PYG{n}{info}\PYG{p}{[}\PYG{l+m+mi}{1}\PYG{p}{]}\PYG{p}{,} \PYG{n}{info}\PYG{p}{[}\PYG{l+m+mi}{2}\PYG{p}{]}\PYG{p}{,} \PYG{n}{info}\PYG{p}{[}\PYG{l+m+mi}{3}\PYG{p}{]}\PYG{p}{,} \PYG{n}{binUmax}\PYG{p}{,}
\PYG{g+gp}{... }                             \PYG{n}{binUmin}\PYG{p}{,} \PYG{n}{Nbin}\PYG{p}{,} \PYG{n}{Mg\PYGZus{}min}\PYG{p}{,} \PYG{n}{outf} \PYG{o}{=} \PYG{l+s+s1}{\PYGZsq{}}\PYG{l+s+s1}{bin\PYGZus{}info\PYGZus{}7200}\PYG{l+s+s1}{\PYGZsq{}}\PYG{p}{)}
\PYG{g+go}{\PYGZlt{}\PYGZlt{}\PYGZlt{} Bin info \PYGZgt{}\PYGZgt{}\PYGZgt{}}
\PYG{g+go}{dU   : 1.07}
\PYG{g+go}{Umax : 250}
\PYG{g+go}{FWHM : 23 degrees}
\PYG{g+go}{f    : 0.6}
\PYG{g+go}{Shape: (457, 457, 100)}
\PYG{g+go}{data : 22.235}
\PYG{g+go}{Saved: bin\PYGZus{}info\PYGZus{}7200.npz}
\end{sphinxVerbatim}

\clearpage

\end{fulllineitems}

\index{grid() (in module grid)@\spxentry{grid()}\spxextra{in module grid}}

\begin{fulllineitems}
\phantomsection\label{\detokenize{grid:grid.grid}}
\pysigstartsignatures
\pysiglinewithargsret{\sphinxcode{\sphinxupquote{grid.}}\sphinxbfcode{\sphinxupquote{grid}}}{\sphinxparam{\DUrole{n}{infile}}\sphinxparamcomma \sphinxparam{\DUrole{n}{n1}}\sphinxparamcomma \sphinxparam{\DUrole{n}{n2}}\sphinxparamcomma \sphinxparam{\DUrole{n}{nrel}}\sphinxparamcomma \sphinxparam{\DUrole{n}{Umax}}\sphinxparamcomma \sphinxparam{\DUrole{n}{FWHM}}\sphinxparamcomma \sphinxparam{\DUrole{n}{f}}\sphinxparamcomma \sphinxparam{\DUrole{n}{Flag}}\sphinxparamcomma \sphinxparam{\DUrole{n}{nstokes}}\sphinxparamcomma \sphinxparam{\DUrole{n}{seed}\DUrole{o}{=}\DUrole{default_value}{None}}}{}
\pysigstopsignatures
\sphinxAtStartPar
This tapers the sky response and grids the \(uv\) plane and computes convolved visibilities, given a fits file.
\begin{equation*}
\begin{split}\begin{align}
\boxed{\mathcal{V}_{cg}^{P}(\vec{\textbf{U}}_g,\nu_a)=\sum_{i}\  \tilde{w}(\vec{\textbf{U}}_g - \vec{\textbf{U}}_i)\ \mathcal{V}^{P}(\vec{\textbf{U}}_i,\nu_a)\ F_i^{P}(\nu_a)}
\end{align}\end{split}
\end{equation*}\begin{quote}\begin{description}
\sphinxlineitem{Parameters}\begin{itemize}
\item {} 
\sphinxAtStartPar
\sphinxstyleliteralstrong{\sphinxupquote{infile}} (\sphinxstyleliteralemphasis{\sphinxupquote{Fits object}}) \textendash{} Fits file  which contains all the data.

\item {} 
\sphinxAtStartPar
\sphinxstyleliteralstrong{\sphinxupquote{n1}} (\sphinxstyleliteralemphasis{\sphinxupquote{int}}) \textendash{} Starting frequency channel number.

\item {} 
\sphinxAtStartPar
\sphinxstyleliteralstrong{\sphinxupquote{n2}} (\sphinxstyleliteralemphasis{\sphinxupquote{int}}) \textendash{} End frequency channel number.

\item {} 
\sphinxAtStartPar
\sphinxstyleliteralstrong{\sphinxupquote{nrel}} (\sphinxstyleliteralemphasis{\sphinxupquote{int}}) \textendash{} \begin{enumerate}
\sphinxsetlistlabels{\arabic}{enumi}{enumii}{}{.}%
\item {} 
\sphinxAtStartPar
If positive then it copies that frequency channel data to all the frequency channels. Flagging Information is same as original fits file.

\item {} 
\sphinxAtStartPar
If \sphinxhyphen{}1 then leaves the data as is.

\item {} 
\sphinxAtStartPar
If \sphinxhyphen{}2 then simulates random gaussian noise visibilities for the baselines with mean \(\mu = 0\) and std \(\sigma = 1\).

\end{enumerate}


\item {} 
\sphinxAtStartPar
\sphinxstyleliteralstrong{\sphinxupquote{Umax}} (\sphinxstyleliteralemphasis{\sphinxupquote{int}}\sphinxstyleliteralemphasis{\sphinxupquote{ or }}\sphinxstyleliteralemphasis{\sphinxupquote{float}}) \textendash{} The maximum value of baseline \(|\vec{\textbf{U}}|_{max}\). Baselines greater than this are rejected.

\item {} 
\sphinxAtStartPar
\sphinxstyleliteralstrong{\sphinxupquote{FWHM}} (\sphinxstyleliteralemphasis{\sphinxupquote{int}}\sphinxstyleliteralemphasis{\sphinxupquote{ or }}\sphinxstyleliteralemphasis{\sphinxupquote{float}}) \textendash{} The full width half maximum (FWHM) of telescope’s primary beam pattern.(in degrees)

\item {} 
\sphinxAtStartPar
\sphinxstyleliteralstrong{\sphinxupquote{f}} (\sphinxstyleliteralemphasis{\sphinxupquote{int}}\sphinxstyleliteralemphasis{\sphinxupquote{ or }}\sphinxstyleliteralemphasis{\sphinxupquote{float}}) \textendash{} Tapering parameter value.

\item {} 
\sphinxAtStartPar
\sphinxstyleliteralstrong{\sphinxupquote{Flag}} (\sphinxstyleliteralemphasis{\sphinxupquote{boolean type}}\sphinxstyleliteralemphasis{\sphinxupquote{ (}}\sphinxstyleliteralemphasis{\sphinxupquote{True}}\sphinxstyleliteralemphasis{\sphinxupquote{ or }}\sphinxstyleliteralemphasis{\sphinxupquote{False}}\sphinxstyleliteralemphasis{\sphinxupquote{)}}) \textendash{} If True it flags the visibility data(puts the data to zero), else does nothing.

\item {} 
\sphinxAtStartPar
\sphinxstyleliteralstrong{\sphinxupquote{nstokes}} (\sphinxstyleliteralemphasis{\sphinxupquote{int}}\sphinxstyleliteralemphasis{\sphinxupquote{ or }}\sphinxstyleliteralemphasis{\sphinxupquote{list}}) \textendash{} The int or the list elements must be within the range of polarizations in the original fits file.

\item {} 
\sphinxAtStartPar
\sphinxstyleliteralstrong{\sphinxupquote{seed}} (\sphinxstyleliteralemphasis{\sphinxupquote{int}}\sphinxstyleliteralemphasis{\sphinxupquote{ or }}\sphinxstyleliteralemphasis{\sphinxupquote{None}}) \textendash{} Default set None, seed value for noise only simulations.

\end{itemize}

\sphinxlineitem{Returns}
\sphinxAtStartPar
The tuple’s first element is a 4D convolved gridded visibility array, where first dimension is polarization axis, next two dimensions are values of convolved visibility at each grid points (square grid) in \(uv\) plane and last one is along frequency axis.
Next element is in the form : (dU, Umax, FWHM, f).

\sphinxlineitem{Return type}
\sphinxAtStartPar
Tuple

\end{description}\end{quote}
\subsubsection*{Examples}

\begin{sphinxVerbatim}[commandchars=\\\{\}]
\PYG{g+gp}{\PYGZgt{}\PYGZgt{}\PYGZgt{} }\PYG{k+kn}{import} \PYG{n+nn}{grid} \PYG{k}{as} \PYG{n+nn}{gd}
\PYG{g+go}{Imported Grid, Import Time : 19/05/26 | 07:28:16 PM}
\PYG{g+gp}{\PYGZgt{}\PYGZgt{}\PYGZgt{} }\PYG{c+c1}{\PYGZsh{} Set the gridding parameters in the TGE code}
\PYG{g+gp}{\PYGZgt{}\PYGZgt{}\PYGZgt{} }\PYG{n}{n1}       \PYG{o}{=} \PYG{l+m+mi}{0}     \PYG{c+c1}{\PYGZsh{} Starting channel number}
\PYG{g+gp}{\PYGZgt{}\PYGZgt{}\PYGZgt{} }\PYG{n}{n2}       \PYG{o}{=} \PYG{l+m+mi}{767}   \PYG{c+c1}{\PYGZsh{} End channel number}
\PYG{g+gp}{\PYGZgt{}\PYGZgt{}\PYGZgt{} }\PYG{n}{Umax}     \PYG{o}{=} \PYG{l+m+mi}{250}   \PYG{c+c1}{\PYGZsh{} Baselines greater than this are rejected}
\PYG{g+gp}{\PYGZgt{}\PYGZgt{}\PYGZgt{} }\PYG{n}{FWHM}     \PYG{o}{=} \PYG{l+m+mi}{23}    \PYG{c+c1}{\PYGZsh{} Primary beam FWHM of the telescope in degrees}
\PYG{g+gp}{\PYGZgt{}\PYGZgt{}\PYGZgt{} }\PYG{n}{f}        \PYG{o}{=} \PYG{l+m+mf}{0.6}   \PYG{c+c1}{\PYGZsh{} Tapering parameter value}
\PYG{g+gp}{\PYGZgt{}\PYGZgt{}\PYGZgt{} }\PYG{n}{Flag}     \PYG{o}{=} \PYG{k+kc}{True}  \PYG{c+c1}{\PYGZsh{} Apply actual flagging of data}
\PYG{g+gp}{\PYGZgt{}\PYGZgt{}\PYGZgt{} }\PYG{n}{nstokes}  \PYG{o}{=} \PYG{p}{[}\PYG{l+m+mi}{0}\PYG{p}{,}\PYG{l+m+mi}{1}\PYG{p}{]} \PYG{c+c1}{\PYGZsh{} 0 for XX and 1 for YY (which polarizations you want to grid)}
\PYG{g+gp}{\PYGZgt{}\PYGZgt{}\PYGZgt{} }\PYG{n}{infile}   \PYG{o}{=} \PYG{l+s+s1}{\PYGZsq{}}\PYG{l+s+s1}{/home/MWA/data\PYGZus{}cutoff/1000000000.fits}\PYG{l+s+s1}{\PYGZsq{}} \PYG{c+c1}{\PYGZsh{} fits file name}
\PYG{g+gp}{\PYGZgt{}\PYGZgt{}\PYGZgt{} }\PYG{n}{nrel}     \PYG{o}{=} \PYG{o}{\PYGZhy{}}\PYG{l+m+mi}{1}  \PYG{c+c1}{\PYGZsh{} data file}
\PYG{g+gp}{\PYGZgt{}\PYGZgt{}\PYGZgt{} }\PYG{n}{GV}\PYG{p}{,} \PYG{n}{info} \PYG{o}{=} \PYG{n}{gd}\PYG{o}{.}\PYG{n}{grid}\PYG{p}{(}\PYG{n}{infile}\PYG{p}{,} \PYG{n}{n1}\PYG{p}{,} \PYG{n}{n2}\PYG{p}{,} \PYG{n}{nrel}\PYG{p}{,} \PYG{n}{Umax}\PYG{p}{,} \PYG{n}{FWHM}\PYG{p}{,} \PYG{n}{f}\PYG{p}{,} \PYG{n}{Flag}\PYG{p}{,} \PYG{n}{nstokes}\PYG{p}{)}
\PYG{g+go}{\PYGZlt{}\PYGZlt{}\PYGZlt{}\PYGZlt{}\PYGZlt{}\PYGZlt{}\PYGZlt{}\PYGZlt{}\PYGZlt{}\PYGZlt{}\PYGZlt{} TGE \PYGZgt{}\PYGZgt{}\PYGZgt{}\PYGZgt{}\PYGZgt{}\PYGZgt{}\PYGZgt{}\PYGZgt{}\PYGZgt{}\PYGZgt{}\PYGZgt{}\PYGZgt{}}
\PYG{g+go}{\PYGZlt{}\PYGZlt{} 19/05/26 | 07:28:16 PM \PYGZgt{}\PYGZgt{}}
\PYG{g+go}{Type  : Data}
\PYG{g+go}{\PYGZlt{}\PYGZlt{} 19/05/26 | 07:30:05 PM \PYGZgt{}\PYGZgt{}}
\PYG{g+go}{Elapsed   :   00:01:49.04}
\PYG{g+go}{\PYGZlt{}\PYGZlt{}\PYGZlt{}\PYGZlt{}\PYGZlt{}\PYGZlt{}\PYGZlt{}\PYGZlt{}\PYGZlt{}\PYGZlt{}\PYGZlt{} Done \PYGZgt{}\PYGZgt{}\PYGZgt{}\PYGZgt{}\PYGZgt{}\PYGZgt{}\PYGZgt{}\PYGZgt{}\PYGZgt{}\PYGZgt{}\PYGZgt{}}
\PYG{g+gp}{\PYGZgt{}\PYGZgt{}\PYGZgt{} }\PYG{n+nb}{print}\PYG{p}{(}\PYG{l+s+sa}{f}\PYG{l+s+s2}{\PYGZdq{}}\PYG{l+s+si}{\PYGZob{}}\PYG{n}{GV}\PYG{o}{.}\PYG{n}{shape}\PYG{+w}{ }\PYG{l+s+si}{= \PYGZcb{}}\PYG{l+s+s2}{\PYGZdq{}}\PYG{p}{)}
\PYG{g+go}{GV.shape = (2, 457, 457, 768)}
\PYG{g+gp}{\PYGZgt{}\PYGZgt{}\PYGZgt{} }\PYG{n+nb}{print}\PYG{p}{(}\PYG{l+s+sa}{f}\PYG{l+s+s2}{\PYGZdq{}}\PYG{l+s+s2}{dU   : }\PYG{l+s+si}{\PYGZob{}}\PYG{n}{info}\PYG{p}{[}\PYG{l+m+mi}{0}\PYG{p}{]}\PYG{l+s+si}{:}\PYG{l+s+s2}{.2f}\PYG{l+s+si}{\PYGZcb{}}\PYG{l+s+se}{\PYGZbs{}}
\PYG{g+gp}{... }\PYG{l+s+s2}{      }\PYG{l+s+se}{\PYGZbs{}n}\PYG{l+s+s2}{Umax : }\PYG{l+s+si}{\PYGZob{}}\PYG{n}{info}\PYG{p}{[}\PYG{l+m+mi}{1}\PYG{p}{]}\PYG{l+s+si}{\PYGZcb{}}\PYG{l+s+se}{\PYGZbs{}}
\PYG{g+gp}{... }\PYG{l+s+s2}{      }\PYG{l+s+se}{\PYGZbs{}n}\PYG{l+s+s2}{FWHM : }\PYG{l+s+si}{\PYGZob{}}\PYG{n}{info}\PYG{p}{[}\PYG{l+m+mi}{2}\PYG{p}{]}\PYG{l+s+si}{\PYGZcb{}}\PYG{l+s+s2}{ degrees}\PYG{l+s+se}{\PYGZbs{}}
\PYG{g+gp}{... }\PYG{l+s+s2}{      }\PYG{l+s+se}{\PYGZbs{}n}\PYG{l+s+s2}{f    : }\PYG{l+s+si}{\PYGZob{}}\PYG{n}{info}\PYG{p}{[}\PYG{l+m+mi}{3}\PYG{p}{]}\PYG{l+s+si}{\PYGZcb{}}\PYG{l+s+s2}{\PYGZdq{}}\PYG{p}{)}
\PYG{g+go}{dU   : 1.07}
\PYG{g+go}{Umax : 250}
\PYG{g+go}{FWHM : 23 degrees}
\PYG{g+go}{f    : 0.6}
\end{sphinxVerbatim}

\end{fulllineitems}




\renewcommand{\indexname}{Index}
\printindex
\end{document}